\chapter{Troubleshooting \LaTeX{} Systematically}
\label{chap:troubleshooting}
\index{troubleshooting}\index{error message}\index{minimal working example}\index{auxiliary files}

The most intimidating \LaTeX{} log often begins with one small structural mistake.
A useful debugging habit is therefore conservative: preserve the failing case,
read the first meaningful error, reduce the search area, test one hypothesis,
and compile again.

\begin{learningoutcomes}
  \item Read the first meaningful error and its source context.
  \item Diagnose braces, environments, commands, files, packages, and Unicode.
  \item Create a minimal working example and isolate a failing section.
  \item Repair citations, references, bibliography builds, floats, and boxes.
  \item Decide when stale auxiliary files should be cleared.
  \item Document a reproducible bug report without sharing confidential content.
\end{learningoutcomes}

\section{The eight-step debugging loop}

Use this sequence before making broad changes:

\begin{enumerate}
  \item Save the source and reproduce the failure.
  \item Read the \emph{first meaningful} error, not only the final stop message.
  \item Note the named file and line, but inspect several lines around it.
  \item Check braces, brackets, math delimiters, and environment pairs nearby.
  \item Reduce the document by commenting out large units or using a minimal file.
  \item Check the class or package documentation for the command and options.
  \item Clear generated auxiliary files only when evidence suggests they are stale.
  \item Rebuild in the required sequence and verify the correction in the PDF.
\end{enumerate}

TeX reports where it noticed a contradiction, which can be later than where the
mistake began. An omitted closing brace in a caption may trigger an error at the
next section. Read backwards through the logical unit.

\begin{tryit}
Open a successful build log and locate the engine version, loaded class, output
PDF line, and any warning. Learning the shape of a healthy log makes unusual
messages easier to recognise.
\end{tryit}

\section{Preserve the evidence}

Before debugging, make a recoverable snapshot. Record the exact command, main
file, engine, and whether the failure appears in a clean build. Do not begin by
updating every package or deleting the entire project; that changes multiple
variables and may destroy useful evidence.

The terminal command

\begin{terminalcode}{Capture a repeatable build}
latexmk -pdf -interaction=nonstopmode -halt-on-error main.tex
\end{terminalcode}

stops at a real error while retaining a log. An editor's summary may shorten or
reorder messages; consult the full \file{.log} when the diagnosis is unclear.

\section{Read an error record}

A typical record contains an error beginning with an exclamation mark, context,
and a line number:

\begin{terminalcode}{Representative undefined-command report}
! Undefined control sequence.
l.42 The result was \qyt{3.2}{\metre}.
\end{terminalcode}

The engine does not know \cmd{qyt}. The likely repair is the intended
\cmd{qty} command from \pkg{siunitx}, assuming that package is loaded. Do not
define a new \cmd{qyt} merely to silence a typographical error.

\begin{pausepredict}
If the line shown by an error contains only \cmd{end}\required{table}, which
earlier elements inside the table would you inspect first? Why can the closing
line be innocent?
\end{pausepredict}

\section{Minimal working examples}

A \term{minimal working example} (MWE) is the smallest complete document that
still demonstrates the problem. Start from a copy, never dismantle the only
working research source.

\begin{latexcode}{Minimal figure test}
\documentclass{article}
\usepackage{graphicx}
\begin{document}
\includegraphics[width=.5\linewidth]{figures/sample.pdf}
\end{document}
\end{latexcode}

Remove unrelated packages and prose while preserving the failure. If removing
a package makes the failure disappear, restore it and test loading order or a
smaller interaction. An MWE answers whether the cause belongs to the content,
the preamble, a file, or the larger project.

When material is confidential, replace names and values while retaining the
same commands and structure. Never post protected data in a public support
forum.

\section{Binary isolation}

For a long file, comment out roughly half the content. If the failure remains,
search that half; otherwise search the removed half. Repeating this process
shrinks the suspect region quickly. Preserve necessary structure: commenting
out the beginning of an environment but not its end creates a new error.

For multi-file books, temporarily use \cmd{includeonly} or comment complete
\cmd{input} lines. Once repaired, restore the complete project and rebuild so
cross-file references and lists are current.

\section{\LaTeX{} First-Aid Clinic}

The following short cases pair a symptom with a local repair. The broken code is
shown as evidence; do not paste it into a working manuscript unchanged.

\subsection{Undefined control sequence}

\begin{latexcode}{Broken: misspelt command}
The sample was \emhp{sterile}.
\end{latexcode}

\begin{latexcode}{Corrected command}
The sample was \emph{sterile}.
\end{latexcode}

If the spelling is correct, check whether the defining package or custom-command
file is loaded before use.

\subsection{Missing closing brace}

\begin{latexcode}{Broken: open argument}
\textbf{Important result
The next sentence is ordinary prose.
\end{latexcode}

\begin{latexcode}{Corrected group}
\textbf{Important result}
The next sentence is ordinary prose.
\end{latexcode}

A missing brace can absorb many later tokens, so the reported line may be far
from the opening command. Editor brace matching helps, but inspect nested
commands manually.

\subsection{Mismatched environment}

\begin{latexcode}{Broken: environment names disagree}
\begin{enumerate}
  \item Calibrate the instrument.
\end{itemize}
\end{latexcode}

\begin{latexcode}{Corrected environment pair}
\begin{enumerate}
  \item Calibrate the instrument.
\end{enumerate}
\end{latexcode}

Every \cmd{begin}\required{name} needs a later
\cmd{end}\required{name} of the same type, properly nested.

\subsection{Math delimiter not closed}

\begin{latexcode}{Broken: unclosed inline mathematics}
The estimate $\hat{\theta} is unbiased.
\end{latexcode}

\begin{latexcode}{Corrected inline mathematics}
The estimate $\hat{\theta}$ is unbiased.
\end{latexcode}

Do not repair this by adding a random dollar sign at the end of the document.
Find the intended boundary between mathematics and prose.

\subsection{File not found}

\begin{latexcode}{Broken on a case-sensitive system}
\includegraphics{Figures/Result.PDF}
\end{latexcode}

If the actual file is \file{figures/result.pdf}, correct the path and case:

\begin{latexcode}{Corrected relative file name}
\includegraphics{figures/result.pdf}
\end{latexcode}

Also verify the main build directory and accepted graphic formats. Avoid
fixing a relative-path problem with an absolute personal path.

\subsection{Missing package or class}

``File \file{name.sty} not found'' means the TeX search path cannot locate a
required package. Confirm that the name is genuine and required, then install it
through the TeX distribution's supported manager or use a documented alternative.
For a publisher class, include the authorised file when redistribution permits.
Do not download an arbitrary similarly named file from an unverified mirror.

\subsection{Unicode character error}

With a current UTF-8 workflow, many characters work directly, but the engine,
font, and context still matter. A copied Unicode minus sign, non-breaking space,
or mathematical symbol may trigger an error or render inconsistently.

Inspect the exact code point. Replace semantic mathematical symbols with \LaTeX{}
commands, for example \cmd{minus} is not a standard command but a normal binary
minus is typed as \code{-} in math mode. Use an engine and font configuration
that intentionally supports needed scripts; do not delete a researcher's name
or linguistic data to make an error disappear.

\subsection{Undefined reference}

\begin{latexcode}{Reference and target}
See \cref{fig:growth-curve}.
% Later:
\label{fig:growth-curve}
\end{latexcode}

Check exact label spelling, uniqueness, and placement after the caption. Then
compile again. Cross-references commonly require two \LaTeX{} passes because the
label value is written to and read from an auxiliary file.

\subsection{Undefined citation or empty bibliography}

Check four layers:

\begin{enumerate}
  \item the citation key exactly matches an entry in the \file{.bib} file;
  \item the database is registered with the correct command and path;
  \item the chosen backend, Biber or BibTeX, actually ran;
  \item \LaTeX{} ran again after the backend created bibliography data.
\end{enumerate}

For a \pkg{biblatex} and Biber workflow, \pkg{latexmk} normally determines the
sequence. An empty bibliography can also mean that no entries were cited and
the style prints only cited works.

\subsection{Package option clash}

An option clash often means a class or package already loaded the same package
with different options. Search the log for the first loading location. Pass
compatible options before the load or remove the duplicate call. Do not load
the same package repeatedly in different preamble fragments.

\section{Errors, warnings, and visual problems}

An error prevents a reliable build. A warning signals a condition requiring
judgement. A PDF can also be wrong without either: a transposed table value,
misleading caption, tiny label, or hidden placeholder is not necessarily a TeX
error.

\subsection{Overfull boxes}

An overfull box extends beyond its allotted width; the log reports an amount.
Inspect the PDF at that location. Common causes are long URLs, unbreakable code,
wide equations, and tables.

Repair the content structure: allow URL breaks, enable code wrapping, split an
equation, revise a column, or use a landscape page. Avoid a global
\cmd{sloppy} setting or shrinking the whole document to conceal one object.

\subsection{Underfull boxes}

An underfull box has poor spacing within a constrained line or page. It may be
harmless in a short caption, narrow table cell, or intentionally ragged region.
Inspect rather than mechanically silencing it. Rewording or widening a column
can help; forced line breaks often create more fragile layout.

\subsection{Float placement}

Figures and tables are floats, so \LaTeX{} balances their preferences with page
constraints. If a float travels too far:

\begin{itemize}
  \item ensure it is referenced near its definition;
  \item use a reasonable width and concise caption;
  \item offer placement choices such as \code{[htbp]};
  \item move the float source after the first explanatory paragraph;
  \item reduce a backlog of large consecutive floats.
\end{itemize}

Use \code{[H]} only with a deliberate reason and supporting package. Repeated
forced placement can create blank space and unstable pagination.

\section{When auxiliary files are stale}

Auxiliary files hold cross-references, citations, contents, lists, index data,
and build state. Clear them when the log suggests corrupted or incompatible
state, after a major tool or class change, or when an impossible old label
survives. Use the build tool's clean command in the specific project:

\begin{terminalcode}{Clean generated \LaTeX{} files}
latexmk -C
latexmk -pdf main.tex
\end{terminalcode}

Cleaning is not a first reflex. It costs useful state and does not repair bad
source. Never use a broad deletion command pointed at an uncertain directory.

\section{Compile in the correct sequence}

A project with citations and an index may need \LaTeX{}, bibliography processing,
index processing, and further \LaTeX{} passes. \pkg{latexmk} automates dependency
detection. If working manually, follow the toolchain documented for that project;
Biber and BibTeX are not interchangeable commands.

Repeated-pass warnings after a stable \pkg{latexmk} build may indicate a value
that changes every run, a conflicting output directory, or an auxiliary file
whose path is not writable. Compare the changing files rather than blindly
increasing the maximum number of passes.

\section{Package documentation and bug reports}

Use \code{texdoc package-name} or the package's authoritative documentation to
verify syntax and compatibility. When requesting help, provide:

\begin{itemize}
  \item a minimal working example;
  \item exact build command and complete first error;
  \item engine, distribution, class, and relevant package versions;
  \item expected and actual behaviour;
  \item whether the problem survives a clean build.
\end{itemize}

Remove confidential content and credentials. Screenshots of an editor can help
with visual UI problems, but selectable error text and a compilable example are
more useful for \LaTeX{} diagnosis.

\begin{goodpractice}
Change one factor at a time and write down the result. The smallest successful
repair teaches more than a copied preamble containing twelve unrelated packages.
After repair, add a test, comment, or README note if the same failure could recur.
\end{goodpractice}

\chapterproject

\begin{miniproject}
Copy the running article to a temporary practice directory. Introduce one
misspelt command, one mismatched environment, one missing figure path, and one
bad citation key, one at a time. For each, save the first meaningful error,
predict the cause, make the smallest correction, and verify the PDF. Finish with
a clean complete build and restore the unbroken project.
\end{miniproject}

\chaptersummary

Systematic debugging preserves evidence and reduces uncertainty. Reproduce the
failure, read the first meaningful error, inspect nearby structure, isolate the
case, consult documentation, and test one repair. Errors in commands, braces,
environments, files, packages, Unicode, references, citations, and build order
have distinct evidence. Box and float warnings require visual judgement.
Auxiliary cleaning is a targeted remedy, and a minimal example is the strongest
foundation for useful support.

\chaptercheck
\begin{chapterchecklist}
  \item The exact failure and build command are reproducible.
  \item The first meaningful error, file, and nearby source were inspected.
  \item Braces, delimiters, and environments are balanced.
  \item The problem survives a reduced example or has been isolated to a unit.
  \item Package syntax and build sequence match authoritative documentation.
  \item Warnings were inspected in the PDF, not merely hidden.
  \item The full project builds after the minimal repair.
\end{chapterchecklist}

\chapterexercisestitle
\begin{chapterexercises}
  \item Explain why the last error in a log may be less useful than the first.
  \item Repair three brace, delimiter, or environment mistakes in supplied code.
  \item Reduce a failing figure document to a minimal working example.
  \item Diagnose an undefined citation across the four bibliography layers.
  \item Give two legitimate and two poor responses to an overfull box.
  \item Explain when cleaning auxiliary files is evidence-based.
  \item Write a privacy-safe bug report containing all reproducibility details.
  \item Challenge: use binary isolation to locate a planted error in a long file
        and record the number of compilation tests required.
\end{chapterexercises}
